\subsection{学界在吸波材料方面的努力}
Zhao等\cite{zhao2022a}通过在碳化硅材料表面镀银,形成导电网络通路,提高其电磁损耗;增强界面极化,提高其衰减常数和阻抗匹配, 从而增强其吸波性能。
Liu等\cite{liu2025}通过超快高温烧结的方式制备了三种新型HEP材料，通过调整组分并触发多种电磁损耗机制，该体系的电磁吸收性能得到了显著提升。
Zhou等\cite{zhou2019}通过化学气相渗透法构建了SiC$_f$/Si$_3$N$_4$复合材料,实现了在 宽温度范围(25–600°C)内X和Ku波段对温度不敏感的电磁波吸收。
Lan和Wang等\cite{lan2020}通过在碳化硅晶须组成的轻质多孔骨架中掺入碳纳米线，热还原以实现碳化硅纳米线网络以产生强界面极化和多重反射，提高电导率。
Cai等\cite{cai2024}通过在1500摄氏度、氮气气氛下烧结碳化硅晶须的方式实现了碳化硅/氮化硅复合气凝胶，将Si$_3$N$_4$纳米带引入碳化硅纳米线网络,增强了阻抗匹配、界面极化和偶极极化。
\subsection{学界目前对于纳米纤维在静电纺织中的演化行为的研究}
